\documentclass[11pt,a4paper]{article}
% LTeX: language=fr
% Character encoding and font configuration for standard pdflatex
\usepackage[utf8]{inputenc}
\usepackage[T1]{fontenc}
\usepackage{mathptmx} % Times New Roman clone for a professional academic look
\usepackage{graphicx}
\usepackage[left=2.5cm,right=2.5cm,top=2cm,bottom=2cm]{geometry}
\usepackage{parskip}   % Automatically adds space between paragraphs and removes indentation
\usepackage[colorlinks=true, urlcolor=black, linkcolor=black]{hyperref}
\usepackage{circuitikz}
\usepackage{float}
\usepackage{subcaption}  % pour subtable


\providecommand{\testAngulaire}{
\begin{circuitikz}
\tikzstyle{every node}=[font=\fontsize{10.2pt}{13.3pt}\selectfont]
\draw [dash dot dot] (17.75,14) -- (22.125,14);
\draw  (19.125,10.75) rectangle (20.75,10.25);
\draw  (22.25,14.25) rectangle (23.875,13.75);
\node [font=\fontsize{10.2pt}{13.3pt}\selectfont, inner xsep=0.080cm, inner ysep=0.085cm, rounded corners=0.020cm] at (24.25,14) {Fixe};
\node [font=\fontsize{14.8pt}{19.2pt}\selectfont, inner xsep=0.080cm, inner ysep=0.085cm, rounded corners=0.020cm] at (23.125,14.625) {$UWB_1$};
\node [font=\fontsize{14.8pt}{19.2pt}\selectfont, inner xsep=0.080cm, inner ysep=0.085cm, rounded corners=0.020cm] at (20,9.875) {$UWB_2$};
\draw [short] (23.25,13.25) -- (23.25,13.25);
\draw [short] (22.75,13.75) -- (22.75,13.75);
\draw [dotted] (23.125,13.625) -- (23.125,10);
\draw [-{Stealth[open,scale=1.5]}, dash dot dot] (20.875,10.5) -- (27,10.5);
\draw [short] (23,13.125) -- (23,13.125);
\draw [short] (23,12.75) -- (23,12.75);
\draw [densely dotted] (22.875,11.75) -- (22.875,11.75);
\draw [densely dotted] (23,12.125) -- (23,12.125);
\draw [densely dotted] (23,12.125) -- (23,12.125);
\draw [{Stealth[scale=1.5]}-{Stealth[scale=1.5]}, ] (20,10.875) -- (23,13.625);
\draw [short] (23,12.25) arc[start angle=-74.74, end angle=-164.74, radius=0.7126];
\node [font=\fontsize{10.2pt}{13.3pt}\selectfont, inner xsep=0.080cm, inner ysep=0.085cm, rounded corners=0.000cm] at (22.375,12.125) {$\theta$};
\draw [dash dot dot] (24.625,14) -- (27,14);
\draw [{Stealth[open,scale=1.5]}-, dash dot dot] (17.625,10.5) -- (19,10.5);
\node [font=\fontsize{10.2pt}{13.3pt}\selectfont, inner xsep=0.080cm, inner ysep=0.085cm, rounded corners=0.000cm] at (23.375,10.25) {y};
\draw [short] (21.375,11.125) -- (21.375,11.125);
\draw [short] (20.75,11.625) -- (20.75,11.625);
\draw [line width=1.2pt, short] (19.625,10.75) -- (20.25,10.75);
\draw [line width=1.2pt, short] (22.75,13.75) -- (23.375,13.75);
\end{circuitikz}
}



\begin{document}

\begin{minipage}[c]{0.30\textwidth}
    \includegraphics[width=\linewidth,keepaspectratio]{Image/Logo_INSA.png}
\end{minipage}
\hfill
\begin{minipage}[c]{0.80\textwidth}
    {\LARGE \textbf{Projet S-IoT : Essaim de drone}}\\[2pt]
    {\large Protocole de calibration }
\end{minipage}

\vspace{0.4cm}
\hrule height 1pt
\vspace{0.8cm}

\begin{center}
    {\huge \textbf{Protocole de Test}} \\[5pt]
    {\large\textit{du 13/11/2026}} \\[5pt]
\end{center}

\tableofcontents

\newpage
Le but de ce document est de définir les tests qui seront réalisés durant la première journée de test en extérieur le 13/11/26. 

Les tests a réalisé sont (par ordre de priorité) : 
\begin{itemize}
    \item Étude de l'angle à 3 drones pour établir la meilleure réception (Précision de 5°)    \item Essai des hauteurs de placement du module UWB pour meilleure réception
    \item Essai de range pour ESP-NOW
    \item Suivi d'un drone avec l'UWB/Communication (pour Z)
    \item Essai de correcteur (Plusieurs valeurs/correcteur)
\end{itemize}


\section{Étude de l'angle : }
Cette étude a pour objectif d'étudier des dispositions angulaires des modules  dans une disposition en triangle pour les drones. 




\section{Étude de la hauteur : }

\section{Essai de range ESP-NOW : }

\section{Suivi d'un drone maître : }

\section{Essai des correcteurs}





% --------------------------------------------------
% Bottom Decorative Rule
% --------------------------------------------------
\vspace{0.4cm}
\vfill
\hrule height 1pt

\end{document}